% Source: https://artofproblemsolving.com/community/c3h1927144_math_problems
% Stable identity: aops-topic-1927144-q05
Place the rectangle with corners $(0,0)$ and $(108,72)$. Since $\gcd(108,72)=36$, the diagonal passes through the lattice points $(3k,2k)$ for $k=0,1,\dots,36$ and through no others: a point $(x,y)$ of the diagonal has $y=\dfrac{2}{3}x$, which is an integer exactly when $3\mid x$. These points cut the diagonal into $36$ congruent pieces, each the diagonal of a $3\times2$ block of squares.

Within one $3\times2$ block, the diagonal from $(0,0)$ to $(3,2)$ crosses the interior vertical lines $x=1,2$ and the interior horizontal line $y=1$ at three different points (none of them a lattice point), and each crossing takes it into a new square. So it passes through $1+3=4$ squares (shaded below).

\begin{asy}
usepackage("amsmath");
size(3.6cm);
defaultpen(fontsize(10pt));
for(int i=0;i<=3;++i) draw((i,0)--(i,2),gray(0.5));
for(int j=0;j<=2;++j) draw((0,j)--(3,j),gray(0.5));
fill((0,0)--(1,0)--(1,1)--(0,1)--cycle,gray(0.8));fill((1,0)--(2,0)--(2,1)--(1,1)--cycle,gray(0.8));
fill((1,1)--(2,1)--(2,2)--(1,2)--cycle,gray(0.8));fill((2,1)--(3,1)--(3,2)--(2,2)--cycle,gray(0.8));
for(int i=0;i<=3;++i) draw((i,0)--(i,2),gray(0.5));
for(int j=0;j<=2;++j) draw((0,j)--(3,j),gray(0.5));
draw((0,0)--(3,2),red+linewidth(1));dot((0,0));dot((3,2));
label("$(0,0)$",(0,0),SW,fontsize(8pt));label("$(3,2)$",(3,2),NE,fontsize(8pt));
label("\mbox{}",(min(currentpicture,user=true)+max(currentpicture,user=true))/2,black); // keeps the text after the figure black when LaTeX typesets the labels inline
\end{asy}

The blocks have disjoint interiors, so the total is $36\cdot4=\boxed{144}$. (This agrees with the general count $m+n-\gcd(m,n)=72+108-36$.)
